\documentclass[UTF8,11pt,a4paper]{ctexart}
\usepackage[a4paper,left=1.45cm,right=1.45cm,top=1.55cm,bottom=1.65cm]{geometry}
\usepackage{amsmath,amssymb,bm}
\usepackage{graphicx}
\usepackage{booktabs,tabularx,array,longtable}
\usepackage{enumitem}
\usepackage[dvipsnames,svgnames]{xcolor}
\usepackage{tikz}
\usetikzlibrary{arrows.meta,positioning,calc,fit,shapes.geometric}
\usepackage[most]{tcolorbox}
\usepackage{hyperref}
\usepackage{caption}
\usepackage{float}
\usepackage{fancyhdr}
\usepackage{titlesec}
\usepackage{microtype}
\IfFontExistsTF{Times New Roman}{\setmainfont{Times New Roman}}{\setmainfont{Liberation Serif}}
\IfFontExistsTF{Noto Sans CJK SC}{\setCJKmainfont{Noto Sans CJK SC}}{}
\IfFontExistsTF{Noto Sans CJK SC}{\setCJKsansfont{Noto Sans CJK SC}}{}
\IfFontExistsTF{DejaVu Sans Mono}{\setmonofont{DejaVu Sans Mono}}{\setmonofont{Latin Modern Mono}}
\definecolor{mainblue}{HTML}{0B5CAD}
\definecolor{softblue}{HTML}{EAF4FF}
\definecolor{darktext}{HTML}{1F2937}
\definecolor{softgray}{HTML}{F5F7FA}
\definecolor{accent}{HTML}{F59E0B}
\definecolor{goodgreen}{HTML}{2E7D32}
\definecolor{badred}{HTML}{B91C1C}
\definecolor{purple}{HTML}{6D5BD0}
\hypersetup{colorlinks=true,linkcolor=mainblue,urlcolor=DarkOrange,citecolor=mainblue}
\setlength{\parindent}{2em}
\setlength{\headheight}{14pt}
\setlength{\parskip}{0.22em}
\setlist[itemize]{leftmargin=2em,itemsep=0.23em,topsep=0.25em}
\setlist[enumerate]{leftmargin=2.2em,itemsep=0.23em,topsep=0.25em}
\renewcommand{\arraystretch}{1.2}
\captionsetup{font=small,labelfont=bf}
\titleformat{\section}{\Large\bfseries\sffamily\color{mainblue}}{\thesection}{0.7em}{}
\titleformat{\subsection}{\large\bfseries\sffamily\color{darktext}}{\thesubsection}{0.6em}{}
\titleformat{\subsubsection}{\normalsize\bfseries\sffamily\color{darktext}}{\thesubsubsection}{0.55em}{}
\pagestyle{fancy}
\fancyhf{}
\fancyhead[L]{\small\sffamily Mixture-of-Recursions(MoR)}
\fancyhead[R]{\small\sffamily \leftmark}
\fancyfoot[C]{\small\thepage}
\renewcommand{\headrulewidth}{0.4pt}
\newtcolorbox{infobox}[1]{enhanced,breakable,colback=softblue,colframe=mainblue!70!black,boxrule=0.8pt,arc=2mm,left=1.5mm,right=1.5mm,top=1mm,bottom=1mm,title=\bfseries #1}
\newtcolorbox{keybox}[1]{enhanced,breakable,colback=softgray,colframe=darktext!45,boxrule=0.6pt,arc=2mm,left=1.5mm,right=1.5mm,top=1mm,bottom=1mm,title=\bfseries #1}
\newtcolorbox{notebox}[1]{enhanced,breakable,colback=accent!8,colframe=accent!80!black,boxrule=0.7pt,arc=2mm,left=1.5mm,right=1.5mm,top=1mm,bottom=1mm,title=\bfseries #1}
\newcommand{\term}[1]{\textbf{\textcolor{mainblue}{#1}}}
\begin{document}

\section{论文基本信息}
\begin{center}
\begin{tabularx}{0.94\linewidth}{>{\bfseries}p{0.18\linewidth}X}
\toprule
题目 & Mixture-of-Recursions(MoR)\\
课件日期 & 2025-07-30\\
研究方向 & 递归 Transformer / 动态计算\\
一句话概括 & 将参数共享、token 级递归路由和递归缓存统一起来,让不同 token 根据难度获得不同计算深度。\\
\bottomrule
\end{tabularx}
\end{center}

\begin{infobox}{摘要要点}
将参数共享、token 级递归路由和递归缓存统一起来,让不同 token 根据难度获得不同计算深度。 本说明按“研究问题—核心机制—基础公式—实验阅读—局限与优化”的顺序重新组织，不保留原始材料的封面、目录和页码对应。
\end{infobox}

\section{研究问题与动机}
这项工作的出发点不是简单地替换某个网络层，而是识别现有范式中最影响\term{质量、效率、可靠性或可部署性}的瓶颈。结合论文所讨论的问题，可以将研究动机概括为以下几点：
\begin{itemize}
\item Recursive Transformers 的局限:固定递归深度,KV 缓存冗余
\item MoR 的核心目标:在单一框架中实现参数共享、自适应计算与高效缓存
\item 前沿工作背景与动机 —— 为什么需要 MoR ?
\item 现有效率方案的不足:参数共享与自适应计算难以兼顾
\item 大模型困境:规模增长导致训练部署成本激增
\item 应用价值:在降低成本的同时保持大模型性能
\end{itemize}

\begin{notebox}{理解重点}
阅读这类论文时，应把“问题是否真实存在”“方法是否直接解决该问题”“实验是否排除了其他解释”分开判断。仅凭更大的模型、更多数据或更长训练得到的提升，不能自动视为结构创新带来的收益。
\end{notebox}

\section{核心思想与整体流程}
将参数共享、token 级递归路由和递归缓存统一起来,让不同 token 根据难度获得不同计算深度。 从信息流角度看，其基本流程可以整理为下图。

\begin{figure}[H]
\centering
\begin{tikzpicture}[>=Latex,node distance=0.78cm]
\tikzstyle{procstep}=[draw,rounded corners,fill=softblue,align=center,minimum width=3.05cm,minimum height=0.95cm,text width=2.75cm]
\node[procstep] (a) {共享递归块};
\node[procstep,right=of a] (b) {token 路由};
\node[procstep,right=of b] (c) {按需递归计算};
\node[procstep,right=of c] (d) {高效解码};
\draw[->,thick,mainblue] (a)--(b);
\draw[->,thick,mainblue] (b)--(c);
\draw[->,thick,mainblue] (c)--(d);
\end{tikzpicture}
\caption{核心信息流与处理阶段。}
\end{figure}

\subsection{方法组成与信息流}
\begin{itemize}
\item Expert-choice 路由:每个递归步骤选择Top-k Token 继续计算
\item 四种参数共享策略:Cycle 、Sequence 及其 “ Middle” 变体
\item 支持深度批处理:不同递归阶段的Token 合并计算,提升GPU 利用率
\item Token-choice 路由:一开始为 Token 分配固定递归次数
\item Middle-Cycle 策略:保留首尾层独特参数,中间层循环共享
\item 路由机制:轻量级路由器为每个 Token 分配递归深度
\item 递归块设计:复用共享层堆叠,替代传统固定深度层
\end{itemize}

\begin{keybox}{结构解析}
\begin{itemize}
\item 5 经典工作回顾 —— Transformer 与递归模型基础 ； Transformer 核心结构：自注意力机制与固定深度的层堆叠 ： （Vaswani et al.,
\end{itemize}
\end{keybox}

上述内容以文字方式保留方法结构，避免把整张幻灯片作为独立页面嵌入正文。

\section{数学与机制理解}
本节给出理解该工作的基础数学结构。公式用于解释方法所属范式，并不替代原论文中的完整推导。
\begin{equation}
\operatorname{Attn}(Q,K,V)=\operatorname{softmax}\!\left(\frac{QK^{\mathsf T}}{\sqrt d}+M\right)V
\end{equation}
注意力方法的关键变量是掩码或路由矩阵 $M$。全注意力允许所有位置交互；稀疏、递归或残差注意力通过限制连接、共享计算或选择历史状态降低成本。

\begin{keybox}{从公式回到方法}
应重点观察论文究竟改变了公式中的哪一部分：是输入表示、连接矩阵、条件变量、参数化方式、训练目标，还是求解与调度过程。真正的创新通常可以在这一层被准确定位。
\end{keybox}


\section{实验结果应如何阅读}
\begin{itemize}
\item 等规模对比:MoR 处理的训练Token 更多(如167M 参数处理27B Token , 基线20B )
\item 实验设置:模型规模135M-1.7B ,训练数据为教育类文本(FineWeb- Edu )
\item 核心指标:验证损失(NLL ,越低越好)、少样本准确率(越高越好)、推 理速度
\item 关键结果:相同计算量下,MoR 参数少50\% ,损失更低、准确率更高
\item 推理提速:动态退出+ 深度批处理,吞吐量最高提升2.06 倍
\item —— 实验结果 MoR 的性能与效率优势
\item Recursive Transformers 的局限:固定递归深度,KV 缓存冗余
\end{itemize}

\begin{notebox}{结果解读}
\begin{itemize}
\item 22 —— 实验结果 MoR 的性能与效率优势 ； 实验设置：模型规模135M-1.7B ，训练数据为教育类文本（FineWeb- Edu ） ； 核心指标：验证损失（NLL ，
\end{itemize}
\end{notebox}

该部分以文字概括实验结论与比较条件，避免用整页截图代替分析。
实验部分不应只看单一最好数字。还需要核对模型规模、训练数据、训练步数、采样或解码预算、硬件条件以及是否使用额外蒸馏、检索或后处理。只有比较条件接近时，结果才能真正说明方法本身的优势。

\section{优点、局限与可优化空间}
\subsection{主要优点}
\begin{itemize}
\item \term{问题与方法对应明确}：论文围绕一个可识别的核心瓶颈展开，方法结构能够直接映射到该瓶颈。
\item \term{可复用的设计思想}：即使不完整复现模型，其中的分解、稀疏化、递归、条件控制或系统优化思想也可迁移到相邻任务。
\item \term{实验提供了多角度证据}：实验通常同时展示主结果、效率比较、案例或消融，使结论不完全依赖单一指标。
\end{itemize}

\subsection{局限与进一步优化}
\begin{itemize}
\item 需要进一步区分\term{结构收益}与更大算力、更多数据、不同训练技巧带来的收益，并在等参数、等计算预算下进行严格比较。
\item 对超参数、输入分布和模型规模的敏感性仍应系统分析；在一个数据集上的最优配置不一定能直接迁移。
\item 实际部署还要考虑显存峰值、并发、延迟抖动、失败案例和鲁棒性，而不仅是离线平均指标。
\item 可尝试将该方法与蒸馏、量化、缓存复用、动态路由或更强的表示学习结合，验证不同优化是否互补。
\end{itemize}

\section{结论}
总体而言，将参数共享、token 级递归路由和递归缓存统一起来，让不同 token 根据难度获得不同计算深度。。 进一步需要注意：MoR 的核心贡献:统一参数共享、自适应路由与高效 KV 缓存;对研究者的启示:动态计算与参数共享结合是高效模型的重要方向。

\end{document}
